\scanpage{104}

que les conditions suivantes soient satisfaites

a) $(a_i)_{0\leq i\leq n}$ est \deleted{libre} strictement libre dans $\unsure{E|\mathcal E}$ (i.e. engendre un sous-Module facteur direct de rg $n+1$)

a$'$) $(a'_i)$ est strictement libre dans $E^\vee$.

b) $\langle a_i,a'_j\rangle=0\qquad 0\leq i<j\leq n,\ j\ne i+1$\quad (ou encore $h_i\subset H_j$ pour ces \unsure{$(i,j)$|$i,j$})

c) \deleted{\unsure{$H_j\ H_i$}} $\langle a_i,a'_{i+1}\rangle\in k^*\qquad(0\leq i\leq n-1)$

\[(\Longleftrightarrow h_i\not\subset H_{i+1}\text{ en chaque fibre},\quad0\leq i\leq n-1)\]

(Ces conditions impliquent \deleted{donc} (28) $\langle a_i,a'_j\rangle=0$ pour $0\leq i,j\leq n$, $i,j$ non consécutifs). Pour qu'on ait

\uline{Dém} La nécessité de a) b) \inserted{c)} est claire, \deleted{\unsure{pour la|prouvons}} \inserted{prouvons} celle de a$'$), \deleted{\unsure{pour cela il faut|prouvons} il suffit de prouver une propriété de stabilité des conditions du Th. 1 par dualité. Or considérons $E_n\subset E_{n+1}=E$, qui est un hyperplan, et soit $s'_0$ une}

On suppose donc qu'on a $s_0$ satisfaisant les conditions dites, on pose. \deleted{$a_{-1}=s_0$, et on définit}

\[a_{-1}=s_0\in E,\qquad a'_{n+1}\in E^\vee\]

défini par

\[\operatorname{Ker}a'_{n+1}=k(a_{-1},a_0\ldots,a_{n-1})\]

\[\text{donc}\quad\left\{\begin{aligned}\langle a_i,a'_{n+1}\rangle&=0&&-1\leq i\leq n-1\\ \inserted{c_n\overset{\mathrm{dfn}}{=}}\langle a_n,a'_{n+1}\rangle&\in k^*\end{aligned}\right.\]

\[\begin{aligned}\text{Considérons}\quad &(a_{-1},a_0,\ldots a_n)\in E^{n+2}\\ &(a'_0,a'_1,\ldots,a'_{n+1})\in E^{\unsure{n+2|-+2|\vee\,n+2}}\end{aligned}\]

et la matrice qu'ils forment, c'est une matrice triangulaire, avec sur la diagonale

% archive sheet number: 103
